\section{Yêu cầu chức năng (Functional requirements)}

Chương này chi tiết hóa các nghiệp vụ mà hệ thống phải thực hiện. Các yêu cầu được xây dựng dựa trên sự phân chia nhiệm vụ của nhóm, đặc biệt tập trung vào các chức năng quản trị nội dung và tương tác khách hàng.

\subsection{Quản trị giao diện và Thông tin chung (Task 1)}
Nhóm chức năng này đảm bảo website luôn hiển thị thông tin chính xác và đồng nhất về thương hiệu:
\begin{itemize}
    \item \textbf{Quản lý Header/Footer:} Admin có thể thay đổi Logo, số điện thoại Hotline, địa chỉ Email và các liên kết mạng xã hội của doanh nghiệp ngay trên trang Dashboard.
    \item \textbf{Hiển thị Banner \& Tin nổi bật:} Hệ thống hỗ trợ cập nhật hình ảnh banner và các thông điệp marketing chính trên trang chủ để thu hút khách hàng.
    \item \textbf{Cấu hình SEO cơ bản:} Cho phép Admin tùy chỉnh tiêu đề trang (Title) và mô tả (Meta Description) cho từng mục để tối ưu hóa sự hiện diện trên các công cụ tìm kiếm.
\end{itemize}

\subsection{Tiếp nhận và Quản lý liên hệ (Task 1)}
Quy trình xử lý yêu cầu của khách hàng được thực hiện khép kín nhằm nâng cao chất lượng dịch vụ:
\begin{itemize}
    \item \textbf{Biểu mẫu liên hệ trực tuyến:} Khách hàng có thể gửi yêu cầu tư vấn, báo giá hoặc phản hồi thông qua form tại trang Liên hệ. Dữ liệu được kiểm tra tính hợp lệ (Validation) trước khi lưu vào database.
    \item \textbf{Hệ thống quản lý phản hồi (Inbox):} Admin có thể xem danh sách các yêu cầu liên hệ, phân loại mức độ quan trọng, ghi chú trạng thái đã xử lý và xóa các phản hồi rác.
    \item \textbf{Thông báo tự động:} Sau khi khách hàng gửi form thành công, hệ thống hiển thị thông báo xác nhận và lưu trữ thời gian thực để Admin dễ dàng theo dõi.
\end{itemize}

\subsection{Quản trị nội dung Giới thiệu doanh nghiệp (Task 2)}
Đảm bảo khách hàng nắm bắt được năng lực và văn hóa của công ty:
\begin{itemize}
    \item \textbf{Trình bày nội dung chi tiết:} Hệ thống cho phép cập nhật lịch sử hình thành, sứ mệnh, tầm nhìn và giá trị cốt lõi của doanh nghiệp.
    \item \textbf{Quản lý đội ngũ nhân sự:} Admin có thể thêm mới hoặc cập nhật thông tin về ban lãnh đạo, các phòng ban chuyên môn để tạo niềm tin với đối tác.
    \item \textbf{Hình ảnh doanh nghiệp:} Hỗ trợ tải lên và hiển thị bộ sưu tập ảnh về văn phòng, hoạt động tập thể của công ty.
\end{itemize}

\subsection{Hệ thống Giải đáp thắc mắc - Q\&A (Task 2)}
Tự động hóa việc hỗ trợ khách hàng thông qua bộ câu hỏi thường gặp:
\begin{itemize}
    \item \textbf{Phân loại chủ đề Q\&A:} Các câu hỏi được nhóm theo từng hạng mục (ví dụ: Dịch vụ, Thanh toán, Bảo mật) để khách hàng dễ dàng tra cứu.
    \item \textbf{Quản lý CRUD Q\&A:} Admin có toàn quyền Thêm, Sửa, Xóa các bộ câu hỏi và câu trả lời trong trang quản trị.
    \item \textbf{Tìm kiếm nhanh:} Tích hợp tính năng tìm kiếm theo từ khóa trong mục Q\&A để nâng cao trải nghiệm người dùng.
\end{itemize}

\subsection{Hệ thống quản trị tập trung (Admin Dashboard)}
Tất cả các chức năng quản trị được tích hợp vào một giao diện chuyên nghiệp sử dụng template \textbf{Srtdash}:
\begin{itemize}
    \item \textbf{Xác thực quyền quản trị:} Chỉ người dùng có tài khoản Admin mới được truy cập vào hệ thống quản lý. Mật khẩu được mã hóa an toàn.
    \item \textbf{Thống kê tổng quan:} Dashboard hiển thị các biểu đồ về số lượng liên hệ mới trong ngày, tổng số bài viết giới thiệu và số câu hỏi Q\&A hiện có.
    \item \textbf{Nhật ký hoạt động:} Ghi nhận các thay đổi quan trọng trên hệ thống để đảm bảo tính minh bạch và bảo mật dữ liệu.
\end{itemize}